\section{Readings we withdrew, and what caused each}
\label{app:withdrawn}

This appendix exists because the alternative is worse. Over the course of this work eight readings of
our own measurements turned out not to survive a check, and in every case the reading we had chosen
was the more interesting of the available ones. The last two were found by an external adversarial
review of the submitted manuscript, after the other six had been written up here --- which is the
most useful thing this appendix can tell a reader about how much to trust the rest. Listing them costs us the appearance of a clean
result and buys the reader the only thing that makes the surviving claims worth anything: a record
of how hard they were pushed. The register in \code{docs/known\_defects.md} carries the code-level
detail; what follows is the manuscript-level consequence of each.

\paragraph{1. ``The two estimands disagree'' (withdrawn; cause: a defective threshold search).} The
analysis plan named two co-primary estimands, $s_{95}$ and the recovery AUC. An earlier version of
Section~\ref{sec:res-primary} reported that they disagreed --- $s_{95}$ a tie at every width, the
AUC favouring the network at $d=50$ and the probe at $d\in\{100,200\}$ --- and treated the
disagreement as the finding, on the grounds that a blunt statistic and a fine one seeing different
things is informative.

The cause was in \code{select\_thresholds}. For the \code{global} and \code{per\_feature} policies it
searched a grid of quantiles of the validation scores and returned the best point on that grid,
without ever scoring $\theta_{\mathrm{fixed}}$, the value it was meant to improve on. On these score
distributions the quantile grid could and did miss, so a decoder could be evaluated at a threshold
that loses to the default on the very validation objective the search maximises. It cost the network
up to five sparsity levels at $d=400$. The fix scores $\theta_{\mathrm{fixed}}$ first, adds a uniform
grid to the quantile grid, and asserts the post-condition that the returned threshold is at least as
good as $\theta_{\mathrm{fixed}}$ on validation; \code{tests/test\_threshold\_selection.py}
re-implements the old search and requires it to fail a bimodal case the new one passes. Both sides of
the comparison moved when it was refitted, which is the reason we trust the correction: it was not a
change that could only help one arm.

With the defect fixed the two estimands agree at all four widths and the disagreement was an
artefact. We report the agreement, and this paragraph, rather than quietly substituting one for the
other.

\paragraph{2. ``The ReLU discards affinely decodable information'' (withdrawn; same cause).} A probe
fitted on the preactivation $\Phi b$ once beat the network by a full sparsity level, and we read that
as a measurement of what the nonlinearity throws away --- a claim we liked, because it reframed a
decoder comparison as a statement about representations. It does not survive the same fix. The
pre-ReLU probe's advantage is now \PrimLfourFiftyPreDiff{} at $d=50$
([\PrimLfourFiftyPreCiLow, \PrimLfourFiftyPreCiHigh], \PrimLfourFiftyPreTies{} of \EsevenSeeds{}
seeds tied), \PrimLfourHundredPreDiff{} at $d=100$ and \PrimLfourTwoHundredPreDiff{} at $d=200$, and
at $d=400$ it reverses: the network leads the \emph{pre}-ReLU probe by
\PrimLfourFourHundredPreDiff{} on \PrimLfourFourHundredPreNetWins{} of \ScSeeds{} models. A quarter
of a sparsity level, with most seeds tied and the sign flipping at the largest width, does not
support a claim about what a nonlinearity discards.

What does survive is confined to the untrained arm, where the effect is an order of magnitude larger
and unanimous: the pre-ReLU probe beats the network by \PrimRandFiftyPreDiff{},
\PrimRandHundredPreDiff{}, \PrimRandTwoHundredPreDiff{} and \PrimRandFourHundredPreDiff{} levels at
the four widths, on every one of the \PrimRandFourHundredPreProbeWins{} and \EsevenSeeds{} models per
cell. On an untrained code the thresholded post-ReLU readout is far worse than an affine read of the
preactivation, and training closes almost all of that gap. We state it in that scoped form and note
that it is a comparison across two stages of one pipeline, not a decoder comparison
(Section~\ref{sec:probes}).

\paragraph{3. ``The frozen-code arm is a failed optimisation'' (withdrawn; cause: scoring an arm
under an objective it was not trained on).} In E8 the arm with a frozen random code trains only
$W_{\mathrm{out}}$ and ends at $\Rreadout\approx\EeightFrozenTwoHundredRreadout$, nearly twice the
cross-talk of the best readout of the code it was given. Since $\Phi^{+}$ attains that optimum for
every code, the obvious reading is that gradient descent failed at a convex-looking problem.

We checked it instead of asserting it. Scored under the loss the arm actually trained on, in its own
post-ReLU representation, on fresh draws from a seed no training used, the frozen $W_{\mathrm{out}}$
\emph{beats} $\Phi^{+}$ by a factor of \EeightFrozenTwoHundredTaskGain{} at $d=200$ on
\EeightFrozenTwoHundredTaskWins{} of \EeightSeeds{} models. The cross-talk is bought, not lost, and
the correct reading is about what co-adaptation does to two objectives rather than about optimisation
difficulty. Section~\ref{sec:res-primary} now says that.

\paragraph{4. ``The advantage of training grows with width'' (withdrawn; cause: choosing the summary
after seeing the result).} Both the trained and the untrained frontier grow with width, so at three
widths it was natural to report that training's advantage grows too --- and as a ratio over
$d\in\{50,100,200\}$ it does. The fourth width breaks it. In absolute sparsity levels, the
operationally meaningful unit, the margin is \MarginFiftyGap, \MarginHundredGap,
\MarginTwoHundredGap, \MarginFourHundredGap: it peaks at $d=200$ and falls, by more than the sampling
noise (a two-sample bootstrap over the independently trained arms gives \MarginDrop{} with interval
[\MarginDropCiLow, \MarginDropCiHigh]). As a ratio it falls at every step.

When we first found the fourth width's numbers we reassured ourselves with the ratio, which was the
reading that preserved the claim. That is the same move as picking a summary statistic after seeing
the result, and we record it as such in \code{docs/decision\_log.md} (D20). The manuscript now leads
with the reading that does not survive and states the surviving claim in its narrowed form: training
buys frontier at every width we measured, and the amount it buys stops growing by $d=400$.

\paragraph{5. ``The subset frontier is a justified shortcut'' (withdrawn; cause: a sufficient
condition used as if it were necessary).} Evaluating $\kappa_i$ for every feature costs $F$ linear
programmes, so early campaigns evaluated a subset chosen by Theorem~\ref{thm:arrow} and reported the
minimum over it as an upper bound. The bound is accurate --- it overestimates by
\SubsetFourHundredOver{} on average even at $d=400$ --- but its justification is not: the theorem is
sufficient for a feature to be hard, not necessary, so the subset need not contain the true argmin.
On untrained codes it always does. On trained codes it drifts, and it drifts with width: the median
leverage rank of the true argmin runs \SubsetFiftyRank, \SubsetHundredRank, \SubsetTwoHundredRank,
\SubsetFourHundredRank, and at $d=400$ the subset locates the argmin in \SubsetFourHundredFound{} of
\ScSeeds{} trained models.

The bias is one-sided and asymmetric between the arms, so a margin computed from subset values is
biased in favour of training by an amount growing with the variable under study. Every frontier
number in this paper is therefore the exact minimum over all $F$ features. We would not use the
shortcut past $d=200$ again, and we report the accuracy and the invalidity together because the first
without the second is how a shortcut survives.

\paragraph{6. Three theorem statements, false as written (corrected; cause: a proof establishing less
than its statement claimed).} An external adversarial audit found that
Theorem~\ref{thm:arrow} and Corollary~\ref{cor:failthresh} were stated for all $h_i<1$ while their
shared proof establishes them only for $h_i\le 1/2$, with explicit unit-norm counterexamples at
$d=2$, $F=3$; and that a proposition concluded a feature with $\kappa_i=1$ is separable at $s=1$ when
the frontier statement requires $\kappa_i>s$, so the correct answer is that it is separable at no
positive sparsity. The statements are corrected here and the counterexamples are regression tests
(\code{tests/test\_arrow\_counterexamples.py}).

No measurement was affected: the corollary is only ever applied at $\min_i h_i$, and the largest
$\min_i h_i$ over all \KdThreeModels{} stage A models is \KdThreeMaxLevMin, inside the valid
region. That is why the tests
passed --- they sampled only the region where the statement is true, which is a gap in the tests and
not a false test. An earlier preprint carries the incorrect statements and will be corrected with an
explicit erratum rather than a silent replacement.

\paragraph{7. ``Stage B repeats stage A's protocol exactly, and denser training states buy frontier''
(withdrawn; cause: comparing two cells without checking they had the same shape).} Stage B was
described in the submitted manuscript as an independent replication of two stage A cells plus a
clean one-variable sparsity contrast. Both halves were false. Its $p=0.01$ cells are at
$F=\SbLoadHigh d$, twice stage A's load, so they are not repeats of anything; and the difference in
effect size between the two campaigns, which we attributed to the imprecision of a cell mean over ten
seeds, is a design difference. Worse, the ``cleanest small result in the campaign'' --- that
$\kappa_{\min}$ rises by a third with the training sparsity --- compared an $F=4d$ cell against an
$F=2d$ one. Made at matched load the comparison reverses: \SbMatchPtwoKappa{}
[\SbMatchPtwoKappaLo, \SbMatchPtwoKappaHi] at $p=0.02$ against \SbMatchPoneKappa{}
[\SbMatchPoneKappaLo, \SbMatchPoneKappaHi] at $p=0.01$, per-model ranges disjoint, in the direction
opposite to the claim.

The claim is withdrawn and Section~\ref{sec:stageb} now reports the null. What survives, correctly
described, is a genuine load axis: the loss--geometry separation and the sign structure of the
decoder comparison hold at twice the feature load. The structural cause is worth more than the
correction: neither \code{primary\_comparison.json} nor \code{all\_feature\_frontier.json} recorded
$F$, so the number distinguishing the two designs was absent from the files every table is generated
from, and no check could have caught the error. Both now record it, and Table~\ref{tab:stageb} carries
it as a column.

\paragraph{8. ``The ceiling analysis evaluates the same objective the selection maximised''
(corrected; cause: the seventh instance of the trap in items 1 and 5).} Section~\ref{sec:probes}
scores the network's own map under the campaign's probe-selection objective. The campaign maximises a
normalised trapezoid over the validation recovery curve; the script computed the plain mean over grid
points, and its own documentation asserted that the mean \emph{was} the selection objective. On a
uniform integer grid the two differ by endpoint reweighting of the same order as the smallest gap the
analysis reports.

Recomputed under the correct functional the result survives and strengthens --- the gaps grow, the
count is unchanged --- so nothing in Section~\ref{sec:probes} is withdrawn. Two things move. The
$d=50$ cell is no longer indistinguishable from zero, so it is no longer offered as a consistency
check. And the $L^2$ cells, which under the wrong functional showed a systematic positive gap, sit at
essentially zero, which is the better outcome: the gap is positive exactly where the network leads,
absent where neither decoder recovers anything, and strongly negative where the probe leads. We
record it because a corrected measurement that happens to help is exactly the kind a reader should be
told about.

\paragraph{What the pattern says.} Four of the eight are the same trap in four different code paths:
an operating point or an objective chosen by maximising something, with no check that it is the thing
it claims to be. Two are a reading chosen from among several that the data permitted, selected after
the fact for being the stronger claim. One --- item 7 --- is a comparison made without checking that
the two things compared had the same shape, which is the same failure at the level of experimental
design rather than of code. We mention the count because a reader deciding how much weight to put on
this paper's surviving claims should know both the rate at which we found our own errors and the
direction they pointed. Only the first pattern is one that more tests would have caught; the last was
caught by recording the variable that distinguished the designs, which is now done.
